% Lösungen Einheit 2 von 4 – Parameter aus der Lagebedingung (zusammenbau v0.1)
\einheitenkopf{Einheit 2 von 4 · Parameter aus der Lagebedingung}

\erg{12}{a) Punkt gegen Ebene – $P$ wird in die Gleichung von $E_k$ eingesetzt. \quad b) $2k + 2 - 3 = 5$, also $k = 3$. \quad c) $A$ und $B$ liefern $6 = 6$, dort fällt $k$ heraus; $C$: $2k = 6$, also $k = 3$. \quad d) z. B. $r = 5$ und $s = -2$, denn $2r + 5s = 0$; $r = s = 0$ ergibt keine Ebenengleichung. \quad e) Der gemeinsame Punkt $(3 | 1 | -1)$ liegt in $E$: $1 - 2 \cdot (-1) = c$, also $c = 3$. \quad f) $2a + 3a - a = 4a = 8$, also $a = 2$: der Punkt $(2 | 2 | 2)$.}
\erg{13}{a) Jana hat den falschen Punkt eingesetzt; mit $P$: $2k + 4 + 2 = 8$, also $k = 1$. \quad b) Liegt $P$ in $E_k$, müssen seine Koordinaten die Gleichung erfüllen; eingesetzt bleibt nur $k$ unbekannt.}
